%% =========================================================================
%% Chapter 4 -- Design, following the lecturer's outline 4.1-4.6 (board of 29/10/2026):
%% 4.1 Class Diagram in Analysis (update, describe in words, explain relationships)
%% 4.2 Data Model (explain relations of tables, generate database, copy screen of DB)
%% 4.3 Class Diagram in Design (update attributes, relationships, methods; explain)
%% 4.4 Layers + class + relation (figures, explain)
%% 4.5 Generate code from 4.4 (present code + structure of packages and classes)
%% 4.6 Conclusion
%% Tables in report/generated are produced from tools/model.py by report/gen_tables.py.
%% =========================================================================
\chapter{Design}\label{ch:design}
\providecommand{\schemaPath}{demo/db/ecom\_schema\_sqlite.sql}

Analysis answered \emph{what} the system means; design answers \emph{how} the software realises it (Lecture~05).
This chapter follows the six points of the lecturer's outline. It starts from the analysis class diagram of
Chapter~\ref{ch:analysis}, turns it into a data model and a database, refines it into a design class diagram with
types and methods, places every class in a layer, and finally generates the code skeleton from the layered model.
Each step produces an artefact that the next step uses -- the analysis class diagram (4.1), the data model and
the database (4.2), the design class diagram (4.3), the layered class model (4.4) and the code (4.5) -- and every
artefact is derived from the same specification file \code{tools/model.py}.

%% =========================================================================
\section{Class Diagram in Analysis}\label{sec:4.1}

The analysis class diagram answers which concepts the system must represent (Lecture~04). Following the lecturer's
rule, it contains no controller, service or repository: every class is a business concept and every relationship a
business constraint with its multiplicity.

\subsection{Update}\label{sec:4.1.1}

The starting point is the class model of Test~01, which used nineteen classes for a general shop (books, clothes,
shoes and electronics). Họa. sells paintings, so the model was updated as listed in Table~\ref{tab:changes41}. The
customer, cart, order, shipping and payment classes are kept; the product side is replaced. The result has
\nAnalysisClasses{} classes.

\begin{longtable}{@{}P{2.6cm}P{3.9cm}P{4.4cm}P{4.1cm}@{}}
\caption{Changes of the analysis class diagram compared with Test~01}\label{tab:changes41}\\
\toprule
\textbf{Point} & \textbf{Test 01} & \textbf{Updated model} & \textbf{Reason}\\
\midrule
\endfirsthead
\toprule
\textbf{Point} & \textbf{Test 01} & \textbf{Updated model} & \textbf{Reason}\\
\midrule
\endhead
\small Products & \small \code{Book}, \code{Clothes}, \code{Shoes}, \code{Electronics}, \code{Laptop}, \code{Mobile} & \small Abstract \code{Product} with the subclasses \code{Painting} and \code{Frame} & \small Họa. sells prints of paintings and frames for them\\
\small Artist & \small -- & \small New class \code{Artist}; an artist \emph{paints} 0..* paintings & \small Customers search by painter (``Monet paintings under 100 dollars'')\\
\small Line items & \small \code{CartItem} and \code{OrderItem} each had four links to the product classes with a \{xor\} constraint & \small One link \emph{refers to} \code{Product}, multiplicity 0..* -- 1 & \small A line always points to exactly one product; the model of Lecture~04 is Cart $\rightarrow$ CartItem $\rightarrow$ Product\\
\small Common attributes & \small Identifier, price and stock repeated in each product class & \small \code{productId}, \code{name}, \code{price}, \code{stockQuantity}, \code{imageUrl}, \code{description} in \code{Product} & \small No duplication; image and description are needed by the AI search\\
\small Card data & \small \code{PayCredit} stored \code{cardNumber} and \code{expiryDate} & \small \code{cardLast4} and \code{authorizationCode} & \small The system shall not store full card numbers (NFR05)\\
\small Multiplicities & \small Added only in the data model & \small Written on the analysis diagram; \code{Order} -- \code{Payment} is 1 -- 0..* & \small Lecture~04 asks for multiplicities in analysis; a failed payment can be retried\\
\small Checkout & \small Dependency \code{Cart} $\dashrightarrow$ \code{Order} «create» & \small Kept & \small Checkout is a transformation, not a permanent association (Lecture~04)\\
\bottomrule
\end{longtable}

\landscapefig{c41_analysis}{Analysis class diagram of Họa.}{Updated analysis class diagram of Họa. (\nAnalysisClasses{} classes; abstract classes in italics). Visual Paradigm diagram \vpsrc{c41_analysis}.}{\label{fig:c41_analysis}\label{fig:analysis41}}

\subsection{Describe in Words}\label{sec:4.1.2}

Figure~\ref{fig:c41_analysis} has four groups of classes.

\textbf{Customers.} A \code{Customer} has an identifier, e-mail, phone and date of birth. The name is a value object
\code{FullName} (first, middle and last name), and the customer keeps one or more \code{Address} objects for
delivery. Two kinds of customers specialise \code{Customer}: \code{CustomerVIP} with a VIP level, a discount rate
and reward points, and \code{CustomerNew} with a registration date and an optional welcome voucher.

\textbf{Products.} \code{Product} is abstract: it holds what every item of the shop has, namely identifier, name,
price, stock quantity, image and description. A \code{Painting} is a print of a painting, described by its style,
subject, dominant colour, medium, year of creation and size in centimetres; each painting is by one \code{Artist},
who has a name, a nationality and life years. A \code{Frame} is sold with a painting and has a material, a colour
and a size label.

\textbf{Cart and order.} A \code{Cart} is the temporary shopping state of one customer and contains
\code{CartItem} lines (quantity and the price shown when the item was added). An \code{Order} is a confirmed
transaction with a date, a status and a total; it contains \code{OrderItem} lines whose \code{unitPrice} is the
price frozen at checkout.

\textbf{Delivery and payment.} \code{Shipping} records the carrier, tracking number, fee, status and delivery date
of an order and points to the delivery \code{Address}. \code{Payment} is abstract (identifier, amount, date,
status) with two forms: \code{PayCash}, paid on delivery, and \code{PayCredit}, authorised by the payment service.

Table~\ref{tab:classes41} in Appendix~B, generated from the model, gives the meaning and the attributes of every
class.

\subsection{Explain Relationships}\label{sec:4.1.3}

Five kinds of relationships are used, each with a precise meaning:
\begin{itemize}
\item \textbf{Composition} (filled diamond): the part cannot exist without the whole and is deleted with it --
\code{Customer} and its \code{FullName}, \code{Cart} and its \code{CartItem} lines, \code{Order} and its
\code{OrderItem} lines.
\item \textbf{Aggregation} (hollow diamond): \code{Customer} groups its \code{Address} objects, but an address is
also used by \code{Shipping}, so it is not deleted with the customer.
\item \textbf{Association}: two independent objects that know each other, with a name and a multiplicity at each
end, for example a customer \emph{places} 0..* orders and an artist \emph{paints} 0..* paintings.
\item \textbf{Generalization} (hollow triangle): \code{CustomerVIP} and \code{CustomerNew} specialise
\code{Customer}, \code{PayCash} and \code{PayCredit} specialise \code{Payment}, \code{Painting} and \code{Frame}
specialise \code{Product}.
\item \textbf{Dependency «create»} (dashed arrow): \code{Cart} creates an \code{Order} at checkout.
\end{itemize}

The multiplicity written at one end says how many objects at this end are linked to \emph{one} object at the
other end; Table~\ref{tab:rel41} in Appendix~B explains every relationship of Figure~\ref{fig:c41_analysis} with its
multiplicities and its business meaning.

Two points of Lecture~04 deserve emphasis. First, \code{CartItem} is not \code{OrderItem}: the cart line is
temporary shopping information whose price may change, the order line is historical transaction information whose
\code{unitPrice} never changes after checkout. Second, checkout is a transformation: it turns
$\{\text{CartItem}_1,\dots,\text{CartItem}_n\}$ into $\{\text{OrderItem}_1,\dots,\text{OrderItem}_n\}$ with the
same products and quantities and leaves the cart empty. This is why \code{Cart} and \code{Order} are linked by a
«create» dependency and not by a permanent association.

%% =========================================================================
\section{Data Model -- Database}\label{sec:4.2}

\subsection{Data Model}\label{sec:4.2.1}

The data model (Figure~\ref{fig:c42_orm}) is derived from the analysis diagram in four steps, the same steps as in
Test~01: every class is marked «ORM-Persistable», meaning that its objects are stored as records; every association
receives a name and multiplicities at both ends; every generalization becomes an association \emph{is a} with
multiplicity 1 at the parent and 0..1 at the child, because the relational model has no inheritance; and the
«create» dependency is dropped, because it describes an operation, not data to store.

\landscapefig{c42_orm}{Data model of Họa.}{Data model: classes marked «ORM-Persistable», generalizations turned into \emph{is a} associations 1 -- 0..1. Visual Paradigm diagram \vpsrc{c42_orm}.}{\label{fig:c42_orm}\label{fig:orm}}

\subsection{From the Data Model to Tables}\label{sec:4.2.2}

Table~\ref{tab:mapping42} lists the mapping rules. They are applied mechanically by \code{tools/gen_sql.py}, so the
database cannot drift from the model.

\begin{longtable}{@{}P{3.4cm}P{6.6cm}P{5.0cm}@{}}
\caption{Rules that turn the data model into tables}\label{tab:mapping42}\\
\toprule
\textbf{Case} & \textbf{Rule} & \textbf{Example}\\
\midrule
\endfirsthead
\toprule
\textbf{Case} & \textbf{Rule} & \textbf{Example}\\
\midrule
\endhead
\small Persistent class & \small One table, name in lower case with underscores; \code{Order} becomes \code{orders} because \code{ORDER} is an SQL keyword & \small \code{CartItem} $\rightarrow$ \code{cart_item}\\
\small Identifier & \small Primary key of type \code{varchar(20)} & \small \code{customer_id}, \code{product_id}, \code{order_id}\\
\small Value object 1 -- 1 & \small Embedded as columns of the owner, no separate table & \small \code{FullName} $\rightarrow$ \code{customer.first_name}, \code{middle_name}, \code{last_name}\\
\small Generalization & \small One table per class (class table inheritance): the child table uses the parent key as primary key and foreign key & \small \code{painting.product_id} $\rightarrow$ \code{product}; \code{pay_cash.payment_id} $\rightarrow$ \code{payment}\\
\small Discriminator & \small The parent table says which child a row belongs to & \small \code{product.product_type} (PAINTING, FRAME), \code{payment.method}\\
\small Association 1 -- * & \small Foreign key on the ``many'' side & \small \code{painting.artist_id}, \code{orders.customer_id}\\
\small Association 1 -- 0..1 & \small Foreign key with a \code{UNIQUE} constraint & \small \code{cart.customer_id}, \code{shipping.order_id}\\
\small Composition with lines & \small Composite primary key; deleting the whole deletes the lines (\code{ON DELETE CASCADE}) & \small \code{cart_item(cart_id, product_id)}, \code{order_item(order_id, product_id)}\\
\small Money, quantities, states & \small Money \code{numeric(12,2)} in US dollars; quantities \code{integer} with \code{CHECK}; states as text with \code{CHECK ... IN (...)} & \small \code{order_item.unit_price}, \code{cart_item.quantity > 0}\\
\small Technical table & \small \code{product_feature} stores the feature vector of each painting for the search & \small \code{product_feature(product_id, model_version, feature_vector, updated_at)}\\
\bottomrule
\end{longtable}

\subsection{Explain Relations of Tables}\label{sec:4.2.3}

Figure~\ref{fig:c42_tables} shows the physical model: each table is drawn as a class with the stereotype «table»,
each column with its SQL type and the markers \{PK\}, \{FK\}, \{U\} (unique) and \{NULL\} (nullable), and each line
is one foreign key from the parent table to the child table, labelled with the relationship name and the
foreign-key column. Table~\ref{tab:tables} in Appendix~B lists the \nTables{} tables with their keys and
Table~\ref{tab:fks} explains the \nFks{} foreign keys with their names and multiplicities.

\landscapefig{c42_tables}{Physical database model of Họa.}{Physical database model: \nTables{} tables «table» with columns, keys and the \nFks{} foreign keys. Visual Paradigm diagram \vpsrc{c42_tables}.}{\label{fig:c42_tables}\label{fig:tables42}}


Foreign keys and \code{UNIQUE} constraints enforce the \emph{upper} bound of a multiplicity, but not the lower bound
1 of a relationship such as \code{orders} -- \code{order_item} (1 -- 1..*). The rules ``an order has at least one
line'' and ``a customer has at least one address'' are therefore checked by the application:
\code{CheckoutService.checkout} refuses an empty cart, and \code{Customer.findAddress} refuses an address of another
customer. In the same way, the discriminator \code{product_type} says which child table must hold the details of a
product; the repositories read and write the parent row and the child row of a product together.

\subsection{Generate Data Base}\label{sec:4.2.4}

The database is generated, not written by hand. \code{tools/gen_sql.py} reads \code{model.TABLES} and writes two
files: \texttt{\schemaPath}, with the \code{CREATE TABLE} statements in the order required by the foreign keys, and
\code{demo/db/seed.sql}, with the demo data (paintings and their artists, frames, the two demo customers and their
orders). When the application starts, \code{EcomApplication.main} calls
\code{Database.initialize(schemaSql, seedSql)}, which creates the SQLite file \code{demo/data/ecom.db} the first
time. Listing~\ref{lst:schema} shows three of the generated tables: the parent table \code{product}, the child table
\code{painting} (class table inheritance and the foreign key to \code{artist}) and the composition table
\code{order_item}.

\IfFileExists{generated/schema_excerpt.sql}{%
\lstinputlisting[language=SQL,caption={Excerpt of the generated schema (\texttt{\schemaPath})},label={lst:schema}]{generated/schema_excerpt.sql}}%
{\assetbox{Excerpt of the generated schema}{report/generated/schema_excerpt.sql}}

The schema is checked automatically: \code{tools/check_code_model.py} opens the real database and compares every
table and column with \code{model.TABLES}, and compares the SQL strings of the repositories with the same list
(Section~\ref{sec:4.5.1}).

\subsection{Copy Screen of DB}\label{sec:4.2.5}

Figures~\ref{fig:db_tables} and~\ref{fig:db_rows} are screenshots of the generated SQLite database opened in a
database viewer: the list of tables, and rows of the \code{painting}, \code{artist} and \code{orders} tables.

\begin{figure}[htbp]
\centering
\begin{minipage}[t]{0.485\linewidth}\centering
\demofig{db_tables}{0.30\textheight}
\caption[Tables of the SQLite database]{Tables of the generated SQLite database \texttt{demo/data/ecom.db}. Screenshot \demosrc{db_tables}.}\label{fig:db_tables}\label{fig:dbtables}
\end{minipage}\hfill
\begin{minipage}[t]{0.485\linewidth}\centering
\demofig{db_rows}{0.30\textheight}
\caption[Rows of the SQLite database]{Rows of the demo data in the SQLite database. Screenshot \demosrc{db_rows}.}\label{fig:db_rows}\label{fig:dbrows}
\end{minipage}
\end{figure}

%% =========================================================================
\section{Class Diagram in Design}\label{sec:4.3}

The design class diagram (Figure~\ref{fig:c43_design}) keeps the \nAnalysisClasses{} domain classes of
Figure~\ref{fig:c41_analysis} and answers a new question: \emph{what does each class hold, with which type, and
which behaviour is it responsible for?} Three enumerations are added for the states (\code{OrderStatus},
\code{ShippingStatus}, \code{PaymentStatus}). Controllers, services and repositories are not domain classes; they
are added in Section~\ref{sec:4.4}.

\landscapefig{c43_design}{Design class diagram of the domain}{Design class diagram of the domain: typed private attributes, public methods, directed associations with role names. Visual Paradigm diagram \vpsrc{c43_design}.}{\label{fig:c43_design}\label{fig:design43}}

\subsection{Update Attributes}\label{sec:4.3.1}

Every attribute becomes private and receives a Java type. The analysis types are refined as follows: money
(\code{price}, \code{cartPrice}, \code{unitPrice}, \code{totalAmount}, \code{shippingFee}, \code{amount},
\code{cashReceived}, \code{changeAmount}, \code{discountRate}) becomes \code{BigDecimal}, because binary floating
point cannot represent cents exactly; dates become \code{LocalDate} (a day, such as \code{dateOfBirth}) or
\code{LocalDateTime} (an instant, such as \code{orderDate}); the free text \code{status} becomes one of the three
enumerations, so an impossible state cannot be written; and \code{Address} receives an identifier
\code{addressId}, needed to choose a delivery address at checkout. Table~\ref{tab:types43} summarises the
refinements; every typed attribute is shown in Figure~\ref{fig:c43_design}.

\begin{longtable}{@{}P{2.7cm}P{3.4cm}P{8.9cm}@{}}
\caption{Refinement of attribute types from analysis to design}\label{tab:types43}\\
\toprule
\textbf{Analysis type} & \textbf{Design type} & \textbf{Attributes}\\
\midrule
\endfirsthead
\toprule
\textbf{Analysis type} & \textbf{Design type} & \textbf{Attributes}\\
\midrule
\endhead
\small \code{float} (money) & \small \code{BigDecimal} & \small \code{price}, \code{cartPrice}, \code{unitPrice}, \code{totalAmount}, \code{shippingFee}, \code{amount}, \code{cashReceived}, \code{changeAmount}, \code{discountRate}\\
\small \code{Date} (a day) & \small \code{LocalDate} & \small \code{dateOfBirth}, \code{registeredDate}, \code{deliveredDate}\\
\small \code{Date} (an instant) & \small \code{LocalDateTime} & \small \code{createdDate}, \code{updatedDate}, \code{orderDate}, \code{paymentDate}\\
\small \code{String} (state) & \small \code{OrderStatus}, \code{ShippingStatus}, \code{PaymentStatus} & \small \code{Order.status}, \code{Shipping.status}, \code{Payment.status}\\
\small -- & \small \code{String} & \small new identifier \code{Address.addressId}\\
\small any & \small private (\code{-}) & \small every attribute; read and written through getters and setters\\
\bottomrule
\end{longtable}

\subsection{Update Methods}\label{sec:4.3.2}

Lecture~05 warns against two opposite mistakes: putting all logic in the controller or in one ``god'' service, and
leaving the domain classes as bare data. Họa. keeps the business rules in the domain classes and lets the services
only coordinate. Table~\ref{tab:ops43} in Appendix~B lists every method with its signature; the important ones are:
\begin{itemize}
\item \code{Cart.addItem(product, quantity)} merges a product already in the cart and rejects a quantity of zero
or less; \code{Cart.calculateTotal()} sums \code{CartItem.getSubtotal()}.
\item \code{OrderItem.fromCartItem(item, current)} (static) creates an order line from a cart line at the
\emph{current} catalogue price: this is the checkout transformation of Lecture~04 at the level of one line.
\item \code{Order.calculateTotal(discountRate)} applies the discount of the customer and adds the shipping fee;
\code{Order.cancel()} is allowed only from \code{PLACED} or \code{PROCESSING} (FR11); \code{Order.advanceTo(next)}
accepts only \code{PLACED} $\rightarrow$ \code{PROCESSING} $\rightarrow$ \code{SHIPPED} $\rightarrow$
\code{DELIVERED} (FR13); \code{Order.isOwnedBy(customerId)} is the ownership rule of NFR01.
\item \code{Product.isAvailable(quantity)} and \code{Product.decreaseStock(quantity)} protect the stock;
\code{Product.keywords()} returns the searchable text, extended by \code{Painting} with style, subject, colour,
medium and artist.
\end{itemize}
Polymorphism replaces conditions on a type field: \code{Product.getCategory()} is abstract and answered by
\code{Painting} and \code{Frame}; \code{Payment.process()} is abstract and answered differently by \code{PayCash}
(cash checked on delivery) and \code{PayCredit} (authorisation code required); \code{Customer.getDiscountRate()}
returns zero and is overridden by \code{CustomerVIP}.


\subsection{Update Relationships -- Explain}\label{sec:4.3.3}

In design, an association becomes a reference held by one class, so each association receives a direction
(navigability) and a role name, which will be the name of the field (Table~\ref{tab:nav43} in Appendix~B). The direction follows
the questions the code must answer. A cart must know its \code{owner} to check that the customer of the session may
check it out, so the analysis link ``customer owns cart'' becomes \code{Cart} $\rightarrow$ \code{Customer}; in the
same way an order knows its \code{customer} and a painting knows its \code{artist}. The whole knows its parts in a
composition (\code{Cart.items}, \code{Order.items}), and an order knows its \code{shipping} and \code{payments}. The
generalizations are kept and become Java inheritance, with \code{Product} and \code{Payment} abstract.

The «create» dependency between \code{Cart} and \code{Order} of the analysis diagram is refined in two places: at
the level of one line, \code{OrderItem} uses \code{CartItem} through \code{fromCartItem}; at the level of the whole
process, the application service \code{CheckoutService} creates the order (Section~\ref{sec:4.4}). This follows
Lecture~05: checkout is an application process, so it belongs to an application service.


%% =========================================================================
\section{Layers + Class + Relation}\label{sec:4.4}

\subsection{Layers and Packages}\label{sec:4.4.1}

Following Lecture~05, Họa. has four layers (Figure~\ref{fig:c44_layers}). The \textbf{presentation} layer receives
requests and presents results (browser views, controllers, the web server); the \textbf{application} layer
coordinates use cases (services, including the search services of A06); the \textbf{domain} layer represents the
business concepts and rules; the \textbf{infrastructure} layer stores data and talks to external technical services
(SQLite repositories, the vector index, the Gemma~4 and Whisper clients).

\begin{figure}[htbp]
\centering
\vpfig{c44_layers}{0.40\textheight}
\caption[Layers and packages of Họa.]{The four layers (packages), their sub-packages and the allowed
dependencies. Visual Paradigm diagram \vpsrc{c44_layers}.}
\label{fig:c44_layers}\label{fig:layers44}
\end{figure}

The dashed arrows are the only allowed dependencies: presentation uses application; application uses domain and
infrastructure; infrastructure uses domain to read and write domain objects. There is no arrow from presentation
to infrastructure, so the user interface never touches the database (NFR04), and the domain depends on no other
layer, so the business rules do not depend on SQLite, HTTP or a model. Lecture~05 stresses that a package is not a
layer: a layer is an architectural responsibility, a package is a way to organise code. Here each layer is a root
package, divided further by function (\code{application.service}, \code{application.search}) or by business area
(\code{domain.cart}, \code{domain.order}, \dots). Table~\ref{tab:layers} in Appendix~B lists the packages and their classes.

\subsection{Classes and Relationships per Layer}\label{sec:4.4.2}

Figure~\ref{fig:c44_classes} places the classes in the layers and draws the relationships that cross a layer
boundary; Figure~\ref{fig:c44_classes_backend} completes the application and infrastructure classes. The domain
classes are drawn by name only; their attributes and methods are in Figure~\ref{fig:c43_design}.

\landscapefig{c44_classes}{Classes and relationships per layer}{Classes per layer and the relationships «use», «create» and realization between them. Visual Paradigm diagram \vpsrc{c44_classes}.}{\label{fig:c44_classes}\label{fig:classes44}}

\landscapefig{c44_classes_backend}{Application and infrastructure classes}{Application and infrastructure classes: services, search classes, repository interfaces with their SQLite implementations, and the AI clients behind their interfaces. Visual Paradigm diagram \vpsrc{c44_classes_backend}.}{\label{fig:c44_classes_backend}}

Three kinds of relationships appear:
\begin{itemize}
\item \textbf{«use»} dependencies between layers. Each browser view calls one controller over HTTP (for example
\code{SearchUI} $\rightarrow$ \code{SearchController}, \code{CartView} $\rightarrow$ \code{CartController} and
\code{CheckoutController}); each controller uses one or two services; the services use repository interfaces, the
search classes and the AI interfaces (for example \code{CheckoutService} uses \code{CartRepository},
\code{OrderRepository}, \code{ProductRepository} and \code{CustomerRepository}). Every dependency is a field set by
the constructor.
\item \textbf{«create»}: \code{CheckoutService} creates the \code{Order}; this is the design form of the checkout
transformation.
\item \textbf{Realization}: each repository interface has one SQLite implementation, for example
\code{SqliteOrderRepository} implements \code{OrderRepository} and \code{SqliteCartRepository} implements
\code{CartRepository}; in the same way \code{GemmaClient} implements \code{LlmClient} and \code{WhisperClient}
implements \code{SpeechClient}.
\end{itemize}
The services depend on the interfaces, never on the SQLite or HTTP classes. This is the dependency inversion
principle: the database or the AI provider can be replaced, or replaced by a fake in tests, without changing the
application layer.

\subsection{Explain}\label{sec:4.4.3}

Lecture~05 assigns four responsibilities to four kinds of classes. Table~\ref{tab:resp44} shows how Họa. applies
them to checkout.

\begin{longtable}{@{}P{2.2cm}P{3.6cm}P{3.9cm}P{5.3cm}@{}}
\caption{The four design responsibilities in Họa.}\label{tab:resp44}\\
\toprule
\textbf{Component} & \textbf{Question (Lecture 05)} & \textbf{Class in Họa.} & \textbf{Responsibility in checkout}\\
\midrule
\endfirsthead
\toprule
\textbf{Component} & \textbf{Question (Lecture 05)} & \textbf{Class in Họa.} & \textbf{Responsibility in checkout}\\
\midrule
\endhead
\small Controller & \small How does a request enter? & \small \code{CheckoutController} & \small Takes the customer id from the session, builds a \code{PayCash} or \code{PayCredit} from the payment method, calls the service\\
\small Service & \small How is a process coordinated? & \small \code{CheckoutService} & \small Runs the steps of the checkout in order; the natural place of a database transaction\\
\small Domain & \small What does the business represent? & \small \code{Cart}, \code{Order}, \code{OrderItem}, \code{Payment} & \small Rules: non-empty cart, frozen price, total with discount, payment processing\\
\small Repository & \small How is data persisted? & \small \code{CartRepository}, \code{OrderRepository}, \code{ProductRepository} & \small Read the cart and the catalogue, save the order, decrease stock, clear the cart\\
\bottomrule
\end{longtable}

\textbf{Checkout across the layers.} Figure~\ref{fig:c44_checkout_seq} shows the design sequence. The steps follow
the body of \code{CheckoutService.checkout} in \code{demo/src}:
\begin{enumerate}
\item The customer confirms the order; \code{CheckoutController.checkout(customerId, cartId, addressId,
paymentMethod)} receives the request with the customer id of the session.
\item The controller creates the payment object and calls \code{CheckoutService.checkout(cartId, customerId,
addressId, payment)}.
\item The service reads the cart with \code{CartRepository.findById(cartId)} and refuses it if its owner is not
the customer of the session, or if it is empty.
\item It reads the customer with \code{CustomerRepository.findById} and creates a new \code{Order} with status
\code{PLACED}.
\item For each \code{CartItem} it reads the product again with \code{ProductRepository.findById}, checks
\code{isAvailable}, and adds \code{OrderItem.fromCartItem(item, current)}: the price is frozen from the catalogue,
not trusted from the cart.
\item It creates the \code{Shipping} to \code{customer.findAddress(addressId)}, then calls
\code{order.calculateTotal(customer.getDiscountRate())}.
\item It fills the payment with the total; a card payment is authorised at once through the polymorphic
\code{process()}, cash stays \code{PENDING} until delivery.
\item It saves the order with \code{OrderRepository.save}, decreases the stock of each line with
\code{ProductRepository.decreaseStock} and empties the cart with \code{CartRepository.clear}.
\item The controller returns the order to the browser view, which shows the confirmation.
\end{enumerate}

\begin{figure}[htbp]
\centering
\vpfig{c44_checkout_seq}{0.55\textheight}
\caption[Design sequence of the checkout]{Design sequence of the checkout through controller, service, domain and
repositories. Visual Paradigm diagram \vpsrc{c44_checkout_seq}.}
\label{fig:c44_checkout_seq}\label{fig:seq44}
\end{figure}

\textbf{Why \texttt{checkout(cartId, customerId, addressId, payment)}.} Lecture~05 proposes
\code{CheckoutService.checkout(cartID) : Order}, because the service can then read the current cart, check its
state and control the transaction. Họa. keeps \code{cartId} as the first parameter and adds three parameters that a
real checkout needs: \code{customerId}, taken from the session and never from the request body, so that the
service can check that the cart belongs to the customer (NFR01); \code{addressId}, one of the customer's addresses;
and the \code{payment} chosen by the customer.

\textbf{Technology mapping.} Lecture~05 asks to design first and choose technology later. Table~\ref{tab:techmap} in
Appendix~B records the choice for Họa.: plain Java classes for controllers, services and domain objects, Java
interfaces with JDBC implementations on a SQLite file for the repositories, browser modules for the views, and HTTP
clients for Gemma~4 and the self-hosted Whisper server.


%% =========================================================================
\section{Generate Code from 4.4}\label{sec:4.5}

\subsection{How the Code Is Generated}\label{sec:4.5.1}

The code is generated from the model of Section~\ref{sec:4.4}, not written from memory. The script
\code{tools/gen_java.py} reads \code{tools/model.py}, the same specification from which the Visual Paradigm
diagrams of this chapter are built, and writes one Java file per class into \code{demo/generated/src}. The browser
views are not generated, because they are JavaScript modules. Table~\ref{tab:uml2java} in Appendix~B lists the mapping.


Listing~\ref{lst:gencheckout} shows a generated file,
\code{demo/generated/src/com/ptit/ecom/application/service/CheckoutService.java}: package, imports, fields and
signatures are complete, and every body that the model does not specify throws an ``unsupported'' exception until
it is implemented.

\IfFileExists{../demo/generated/src/com/ptit/ecom/application/service/CheckoutService.java}{%
\lstinputlisting[language=Java,caption={Generated skeleton of \texttt{CheckoutService}},label={lst:gencheckout}]{../demo/generated/src/com/ptit/ecom/application/service/CheckoutService.java}}%
{\assetbox{Generated skeleton}{demo/generated/src/com/ptit/ecom/application/service/CheckoutService.java}}

The implementation in \code{demo/src/main/java} starts from these files and keeps every name: class, method,
parameter and field names, and the table and column names used in SQL, must be exactly those of the model. The
script \code{tools/check_code_model.py} enforces this rule by comparing the source code, the JavaScript views and the
real database with \code{model.py}; its report, the second part of \code{build/test_results.txt}, finds no
difference -- its ``NOTE'' lines only mark private helper fields, which the naming rule allows -- and is reproduced in
Appendix~B (Section~\ref{sec:B.7}).

\subsection{Present Structure of Packages and Classes}\label{sec:4.5.2}

The implemented project contains \javaFileCount{} Java files (\javaLineCount{} lines) and \jsFileCount{} JavaScript
view modules. Its package tree mirrors the layers of Figure~\ref{fig:c44_layers}: one folder per layer, one
sub-folder per package, one file per class.

\txtfileor{generated/code_tree.txt}{../build/code_tree5.txt}

The other folders of \code{demo} hold what the code needs at run time: \code{web} (pages, styles, the view modules
and the painting images), \code{db} (the generated schema and seed), \code{data} (the SQLite database and the
catalogue), \code{config} (\code{app.properties}), \code{voice-server} (the self-hosted Whisper server) and
\code{generated} (the skeleton of Section~\ref{sec:4.5.1}).

\subsection{Present Code}\label{sec:4.5.3}

\code{CheckoutService} is the class whose generated skeleton is shown in Listing~\ref{lst:gencheckout}; the listing
below is the same class after implementation. It is the application service of Lecture~05: its steps are those of
Section~\ref{sec:4.4.3}, every field and signature is unchanged from the skeleton, and the domain methods it calls
(\code{Cart.isEmpty}, \code{OrderItem.fromCartItem}, \code{Order.calculateTotal}) carry the business rules.

\excerpt{Java}{CheckoutService.java}

The other classes follow the same pattern. \code{OrderService} answers ``find my order'' and ``where is my latest
order?'' of A06 by reading the order through \code{OrderRepository.findOrder} and then applying the ownership rule
\code{Order.isOwnedBy}, so an order of another customer is reported as not found (NFR01).
\code{SqliteOrderRepository} implements the repository with JDBC, using only the tables and columns of the model
and parameterised statements, so a value typed by a user can never change a query. Chapter~\ref{ch:code} presents
the code of the search, AI and web classes.

%% =========================================================================
\section{Conclusion}\label{sec:4.6}

The design turns the meaning of Chapter~\ref{ch:analysis} into a structure that can be programmed. The analysis
diagram was updated for the painting shop; the data model and the \nTables{} tables were derived and generated from
it; the design classes received types, methods and directions; every class was placed in one of four layers with
explicit dependencies; and the Java skeleton was generated from that model. Table~\ref{tab:trace46} in Appendix~B traces the main
requirements through the artefacts of this chapter to the code.


Two limits remain. First, the demo uses one SQLite file and runs the checkout steps through separate repository
calls; a production system would wrap them in one database transaction and support concurrent buyers of the same
print. Second, the attribute vector of the image search is deliberately simple (Chapter~\ref{ch:ai}); the table
\code{product_feature} is ready for a learned embedding. Chapter~\ref{ch:code} presents the complete code, the
user interface and the tests.
